\documentclass[11pt,a4paper]{article}
\usepackage[T1]{fontenc}
\usepackage[utf8]{inputenc}
\usepackage{amsmath,amssymb,xcolor}
\usepackage[margin=25mm]{geometry}
\usepackage[colorlinks=true,urlcolor=blue]{hyperref}
\definecolor{nexorablue}{HTML}{2459B5}
\setlength{\parindent}{0pt}
\setlength{\parskip}{7pt}
\setlength{\emergencystretch}{2em}
\renewcommand{\familydefault}{\sfdefault}
\begin{document}

{\small\bfseries\color{nexorablue} NEXORA CPGE}\par
{\small MPSI / NEX-CIR-04 / Corrigé}\par
{\LARGE\bfseries Énergie reçue sous deux impulsions}\par
\vspace{0.25em}\hrule\vspace{0.6em}

\section*{Charge transférée}
Les contributions sont $30\times0{,}20=6\,\mathrm{mC}$ et $-10\times0{,}10=-1\,\mathrm{mC}$. La charge algébrique totale vaut $5\,\mathrm{mC}$.

\section*{Puissances et énergie}
En convention récepteur, les puissances reçues sont $120$ puis $-40\,\mathrm{mW}$. Le dipôle cède de l’énergie pendant le second intervalle. Le bilan vaut $120\times0{,}20-40\times0{,}10=20\,\mathrm{mJ}$.


\end{document}
